% GENERATED FILE. Edit canonical lessons, then run npm run generate:books.
% canonical-source-hash: fnv1a64:bc39ca0dce8616fe
% canonical-lessons: KA-C321-nyayalaya, KA-C321-nyayadhisa, KA-C321-aropisu, KA-C321-tirpu, KA-C321-sikse

\chapter{In Court}
\label{ch:ka-in-court}
\hlchaptermodality{321}

\begin{chapteropening}
\textbf{By the end of this chapter:} \emph{I can say a court of law, a judge, to accuse, a verdict and a punishment.}

Complete the last lesson of chapter 321: I can say a court of law, a judge, to accuse, a verdict and a punishment.
\end{chapteropening}

\section[\texorpdfstring{\kn{ನ್ಯಾಯಾಲಯ} (nyāyālaya)}{ನ್ಯಾಯಾಲಯ (nyāyālaya)}]{\kn{ನ್ಯಾಯಾಲಯ} (nyāyālaya) — a court of law}

\label{lesson:KA-C321-nyayalaya}



\begin{quote}
Before the new one: say the Kannada for to broadcast, then the Kannada for a headline.
\end{quote}

\subsection*{You'll want to know: \kn{ನ್ಯಾಯಾಲಯ}}
\textbf{\kn{ನ್ಯಾಯಾಲಯ}} — \emph{nyāyālaya} — \textquotedblleft{}a court of law\textquotedblright{}.

\begin{quote}
\kn{ಅವನ} \kn{ಭವಿಷ್ಯ} \kn{ನ್ಯಾಯಾಲಯದ} \kn{ನಿರ್ಧಾರದ} \kn{ಮೇಲೆ} \kn{ಅವಲಂಬಿಸಿದೆ}.
\end{quote}

\begin{quote}
\emph{avana bhaviṣya nyāyālayada nirdhārada mēle avalambiside.}
\end{quote}

\begin{quote}
His future depends on the court's decision.
\end{quote}

\kn{ನ್ಯಾಯ}, justice, and \kn{ಆಲಯ}, a house: the house of justice.

\begin{cousinweb}
\textbf{\kn{ನ್ಯಾಯಾಲಯ}} (\emph{nyāyālaya}) is borrowed from Sanskrit, a compound of \emph{nyāya}, \textquotedblleft{}rule, method, a judicial sentence\textquotedblright{} (CDIAL 7614), and \emph{ālaya}, \textquotedblleft{}a house, a dwelling\textquotedblright{} (CDIAL 1366).
\end{cousinweb}

\begin{grammarlens}[title={What you've built}]
One more word for telling what happened.
\end{grammarlens}

\subsection*{Your turn}
\begin{itemize}
\raggedright
  \item \emph{Say it:} \emph{nyāyālaya}
  \item \emph{Say it:} \emph{nyāyālaya}, once more
  \item \emph{Say it:} \emph{avana bhaviṣya nyāyālayada nirdhārada mēle avalambiside.} — His future depends on the court's decision.
  \item \emph{Recall:} say the Kannada for to broadcast, then the Kannada for a headline, then say \emph{nyāyālaya} again
\end{itemize}

\begin{tcolorbox}[breakable,colback=teal!4,colframe=teal!35!black,title={Before you move on}]
What does \kn{ನ್ಯಾಯಾಲಯ} mean? (\textquotedblleft{}A court of law\textquotedblright{}.) Say it once more.
\end{tcolorbox}
\section[\texorpdfstring{\kn{ನ್ಯಾಯಾಧೀಶ} (nyāyādhīśa)}{ನ್ಯಾಯಾಧೀಶ (nyāyādhīśa)}]{\kn{ನ್ಯಾಯಾಧೀಶ} (nyāyādhīśa) — a judge}

\label{lesson:KA-C321-nyayadhisa}



\begin{quote}
Before the new one: say the Kannada for a headline, then the Kannada for a court of law.
\end{quote}

\subsection*{You'll want to know: \kn{ನ್ಯಾಯಾಧೀಶ}}
\textbf{\kn{ನ್ಯಾಯಾಧೀಶ}} — \emph{nyāyādhīśa} — \textquotedblleft{}a judge\textquotedblright{}.

\begin{quote}
\kn{ನ್ಯಾಯಾಧೀಶರು} \kn{ಎರಡೂ} \kn{ಕಡೆಯ} \kn{ಮಾತು} \kn{ಕೇಳಿದರು}.
\end{quote}

\begin{quote}
\emph{nyāyādhīśaru eraḍū kaḍeya mātu kēḷidaru.}
\end{quote}

\begin{quote}
The judge heard both sides.
\end{quote}

Plural: \kn{ನ್ಯಾಯಾಧೀಶರು}, \emph{nyāyādhīśaru}.

People speak of a judge with the respectful plural: \kn{ನ್ಯಾಯಾಧೀಶರು} \kn{ಬಂದರು}.

\begin{cousinweb}
\textbf{\kn{ನ್ಯಾಯಾಧೀಶ}} (\emph{nyāyādhīśa}) is borrowed from Sanskrit, a compound of \emph{nyāya}, \textquotedblleft{}justice\textquotedblright{} (CDIAL 7614), and \emph{adhīśa}, \textquotedblleft{}a lord, a master\textquotedblright{} (CDIAL 267).
\end{cousinweb}

\begin{grammarlens}[title={What you've built}]
One more word for telling what happened.
\end{grammarlens}

\subsection*{Your turn}
\begin{itemize}
\raggedright
  \item \emph{Say it:} \emph{nyāyādhīśa}
  \item \emph{Say it:} \emph{nyāyādhīśa}, once more
  \item \emph{Say it:} \emph{nyāyādhīśaru eraḍū kaḍeya mātu kēḷidaru.} — The judge heard both sides.
  \item \emph{Recall:} say the Kannada for a headline, then the Kannada for a court of law, then say \emph{nyāyādhīśa} again
\end{itemize}

\begin{tcolorbox}[breakable,colback=teal!4,colframe=teal!35!black,title={Before you move on}]
What does \kn{ನ್ಯಾಯಾಧೀಶ} mean? (\textquotedblleft{}A judge\textquotedblright{}.) Say it once more.
\end{tcolorbox}
\section[\texorpdfstring{\kn{ಆರೋಪಿಸು} (ārōpisu)}{ಆರೋಪಿಸು (ārōpisu)}]{\kn{ಆರೋಪಿಸು} (ārōpisu) — to accuse}

\label{lesson:KA-C321-aropisu}



\begin{quote}
Before the new one: say the Kannada for a court of law, then the Kannada for a judge.
\end{quote}

\subsection*{You'll want to know: \kn{ಆರೋಪಿಸು}}
\textbf{\kn{ಆರೋಪಿಸು}} — \emph{ārōpisu} — \textquotedblleft{}to accuse\textquotedblright{}.

\begin{quote}
\kn{ಅಂಗಡಿಯವನು} \kn{ಆ} \kn{ಹುಡುಗ} \kn{ದುಡ್ಡು} \kn{ಕದ್ದ} \kn{ಅಂತ} \kn{ಆರೋಪಿಸಿದನು}.
\end{quote}

\begin{quote}
\emph{aṅgaḍiyavanu ā huḍuga duḍḍu kadda anta ārōpisidanu.}
\end{quote}

\begin{quote}
The shopkeeper accused that boy of stealing money.
\end{quote}

Past (I …): \kn{ಆರೋಪಿಸಿದೆ}, \emph{ārōpiside}.

The charge comes before \kn{ಅಂತ}. The charge itself is an \kn{ಆರೋಪ}, and \kn{ಆರೋಪಿ} is the accused.

\begin{cousinweb}
\textbf{\kn{ಆರೋಪಿಸು}} (\emph{ārōpisu}) is built on Sanskrit \emph{āropa}, \textquotedblleft{}placing on, charging with\textquotedblright{}, the noun beside \emph{āropaṇa}, \textquotedblleft{}placing on\textquotedblright{} (Monier-Williams), with the verb ending -\emph{isu}.
\end{cousinweb}

\begin{grammarlens}[title={What you've built}]
One more word for telling what happened.
\end{grammarlens}

\subsection*{Your turn}
\begin{itemize}
\raggedright
  \item \emph{Say it:} \emph{ārōpisu}
  \item \emph{Say it:} \emph{ārōpisu}, once more
  \item \emph{Say it:} \emph{aṅgaḍiyavanu ā huḍuga duḍḍu kadda anta ārōpisidanu.} — The shopkeeper accused that boy of stealing money.
  \item \emph{Recall:} say the Kannada for a court of law, then the Kannada for a judge, then say \emph{ārōpisu} again
\end{itemize}

\begin{tcolorbox}[breakable,colback=teal!4,colframe=teal!35!black,title={Before you move on}]
What does \kn{ಆರೋಪಿಸು} mean? (\textquotedblleft{}To accuse\textquotedblright{}.) Say it once more.
\end{tcolorbox}
\section[\texorpdfstring{\kn{ತೀರ್ಪು} (tīrpu)}{ತೀರ್ಪು (tīrpu)}]{\kn{ತೀರ್ಪು} (tīrpu) — a verdict; a judgement}

\label{lesson:KA-C321-tirpu}



\begin{quote}
Before the new one: say the Kannada for a judge, then the Kannada for to accuse.
\end{quote}

\subsection*{You'll want to know: \kn{ತೀರ್ಪು}}
\textbf{\kn{ತೀರ್ಪು}} — \emph{tīrpu} — \textquotedblleft{}a verdict; a judgement\textquotedblright{}.

\begin{quote}
\kn{ನ್ಯಾಯಾಧೀಶರು} \kn{ಸೋಮವಾರ} \kn{ತೀರ್ಪು} \kn{ಕೊಟ್ಟರು}.
\end{quote}

\begin{quote}
\emph{nyāyādhīśaru sōmavāra tīrpu koṭṭaru.}
\end{quote}

\begin{quote}
The judge gave the verdict on Monday.
\end{quote}

A verdict is given: \kn{ತೀರ್ಪು} \kn{ಕೊಡು} or, more formally, \kn{ತೀರ್ಪು} \kn{ನೀಡು}.

\begin{cousinweb}
\textbf{\kn{ತೀರ್ಪು}} (\emph{tīrpu}) is native Dravidian: DEDR 3278 lists Kannada \emph{tīrpu}, \textquotedblleft{}finishing, concluding\textquotedblright{}, under \emph{tīr}, \textquotedblleft{}to end, be settled, be decided\textquotedblright{}. A verdict settles a case.
\end{cousinweb}

\begin{grammarlens}[title={What you've built}]
One more word for telling what happened.
\end{grammarlens}

\subsection*{Your turn}
\begin{itemize}
\raggedright
  \item \emph{Say it:} \emph{tīrpu}
  \item \emph{Say it:} \emph{tīrpu}, once more
  \item \emph{Say it:} \emph{nyāyādhīśaru sōmavāra tīrpu koṭṭaru.} — The judge gave the verdict on Monday.
  \item \emph{Recall:} say the Kannada for a judge, then the Kannada for to accuse, then say \emph{tīrpu} again
\end{itemize}

\begin{tcolorbox}[breakable,colback=teal!4,colframe=teal!35!black,title={Before you move on}]
What does \kn{ತೀರ್ಪು} mean? (\textquotedblleft{}A verdict; a judgement\textquotedblright{}.) Say it once more.
\end{tcolorbox}
\section[\texorpdfstring{\kn{ಶಿಕ್ಷೆ} (śikṣe)}{ಶಿಕ್ಷೆ (śikṣe)}]{\kn{ಶಿಕ್ಷೆ} (śikṣe) — a punishment; a sentence}

\label{lesson:KA-C321-sikse}



\begin{quote}
Before the new one: say the Kannada for to accuse, then the Kannada for a verdict.
\end{quote}

\subsection*{You'll want to know: \kn{ಶಿಕ್ಷೆ}}
\textbf{\kn{ಶಿಕ್ಷೆ}} — \emph{śikṣe} — \textquotedblleft{}a punishment; a sentence\textquotedblright{}.

\begin{quote}
\kn{ಕಳ್ಳನಿಗೆ} \kn{ಎರಡು} \kn{ವರ್ಷದ} \kn{ಶಿಕ್ಷೆ} \kn{ಆಯಿತು}.
\end{quote}

\begin{quote}
\emph{kaḷḷanige eraḍu varṣada śikṣe āyitu.}
\end{quote}

\begin{quote}
The thief got a two-year sentence.
\end{quote}

The one punished takes -\kn{ಗೆ}: \kn{ಕಳ್ಳನಿಗೆ} \kn{ಶಿಕ್ಷೆ} \kn{ಆಯಿತು}, the thief was punished.

\begin{cousinweb}
\textbf{\kn{ಶಿಕ್ಷೆ}} (\emph{śikṣe}) is borrowed from Sanskrit \emph{śikṣā}, which meant \textquotedblleft{}learning, study\textquotedblright{} and also \textquotedblleft{}punishment, correction\textquotedblright{} (CDIAL 12432). Kannada kept the second sense.
\end{cousinweb}

\begin{grammarlens}[title={What you've built}]
One more word for telling what happened.
\end{grammarlens}

\subsection*{Your turn}
\begin{itemize}
\raggedright
  \item \emph{Say it:} \emph{śikṣe}
  \item \emph{Say it:} \emph{śikṣe}, once more
  \item \emph{Say it:} \emph{kaḷḷanige eraḍu varṣada śikṣe āyitu.} — The thief got a two-year sentence.
  \item \emph{Recall:} say the Kannada for to accuse, then the Kannada for a verdict, then say \emph{śikṣe} again
\end{itemize}

\begin{tcolorbox}[breakable,colback=teal!4,colframe=teal!35!black,title={Before you move on}]
What does \kn{ಶಿಕ್ಷೆ} mean? (\textquotedblleft{}A punishment; a sentence\textquotedblright{}.) Say it once more.
\end{tcolorbox}
